\documentclass[../BTL.tex]{subfiles}
\begin{document}

% ─────────────────────────────────────────────────────────
\section{Tổng quan về Hệ thống Gợi ý}
\label{section:2.1}

\subsection{Định nghĩa}
\label{section:2.1.1}

Hệ thống gợi ý (\textit{Recommender System} hoặc \textit{Recommendation System}) là
một dạng hệ thống lọc thông tin có nhiệm vụ dự đoán xếp hạng (\textit{rating}) hoặc
sở thích mà một người dùng cụ thể sẽ dành cho một mục (\textit{item}) nhất định, từ
đó đưa ra các gợi ý phù hợp nhất. Khái niệm này được hình thành từ những năm 1990
với hệ thống lọc cộng tác GroupLens~\cite{resnick1994grouplens}, và kể từ đó đã trở
thành một lĩnh vực nghiên cứu sôi động trong cộng đồng học máy và khai phá dữ liệu.

\subsection{Phân loại}
\label{section:2.1.2}

Dựa trên phương pháp tiếp cận, hệ thống gợi ý được chia thành ba nhóm
chính~\cite{ricci2011recommender}:

\begin{itemize}[itemsep=0.4em]
    \item \textbf{Content-based Filtering} (Lọc theo nội dung): Gợi ý các mục có
    nội dung/đặc trưng tương tự với những mục người dùng đã thích trong quá khứ.
    Ưu điểm: không cần dữ liệu từ người dùng khác. Nhược điểm: bị giới hạn trong
    vùng nội dung quen thuộc (\textit{over-specialization}). Phương pháp này
    không được sử dụng trong đồ án.

    \item \textbf{Collaborative Filtering} (Lọc cộng tác): Gợi ý dựa trên
    hành vi/đánh giá của cộng đồng người dùng. Không cần biết nội dung sản phẩm,
    chỉ cần ma trận tương tác \textit{user-item}. Đây là phương pháp được sử dụng
    chủ yếu trong đồ án này.

    \item \textbf{Hybrid Approach} (Kết hợp): Kết hợp cả hai phương pháp trên để
    tận dụng ưu điểm và khắc phục nhược điểm của từng loại. Được ứng dụng bởi các
    hệ thống lớn như Netflix, YouTube.
\end{itemize}

% ─────────────────────────────────────────────────────────
\section{Collaborative Filtering}
\label{section:2.2}

\subsection{Nguyên lý hoạt động}
\label{section:2.2.1}

\textit{Collaborative Filtering} (CF) dựa trên giả định cơ bản: những người dùng có
hành vi tương tự trong quá khứ sẽ có sở thích tương tự trong tương
lai~\cite{sarwar2001item}. CF khai thác ma trận tương tác \textit{user-item} $R$
(kích thước $n \times m$) trong đó $R[u][i]$ là \textit{rating} mà người dùng $u$
dành cho sản phẩm $i$. Các ô chưa có giá trị (NaN) là những \textit{rating} cần
được dự đoán.

CF được chia thành hai nhánh chính: \textbf{Memory-based CF} (dựa trên bộ nhớ, gồm
\textit{User-based} và \textit{Item-based}) và \textbf{Model-based CF} (dựa trên mô
hình, gồm \textit{Matrix Factorization} và các mô hình học sâu). Đồ án này triển
khai cả ba phương pháp: \texttt{UserBasedCF}, \texttt{ItemBasedCF} và
\texttt{MatrixFactorization}, được cài đặt hoàn toàn từ đầu bằng NumPy (không sử
dụng thư viện CF chuyên dụng).

\textbf{Vấn đề Cold Start:} CF gặp khó khăn khi người dùng mới hoặc sản phẩm mới
chưa có đủ dữ liệu tương tác — không thể tính được \textit{similarity} vì vector
\textit{rating} gần như rỗng. Trong hệ thống \textbf{CineRec}, trường hợp này được
xử lý bằng cách \textit{fallback} về \textit{rating} trung bình của người dùng
($\bar{r}_u$) khi không tìm được láng giềng hợp lệ (similarity~$> 0$).

\subsection{User-based Collaborative Filtering}
\label{section:2.2.2}

\textit{User-based CF} tìm kiếm $K$ người dùng có sở thích tương đồng nhất với
người dùng mục tiêu $u$ ($K$ \textit{nearest neighbors}), sau đó tổng hợp
\textit{rating} của họ để dự đoán \textit{rating} cho sản phẩm $i$ chưa được đánh
giá. Trong đồ án này, $K = 30$.

Độ tương đồng giữa hai người dùng được đo bằng \textbf{mean-centered cosine
similarity} — trừ \textit{rating} trung bình của từng người dùng để loại bỏ
\textit{bias} đánh giá (xu hướng cho điểm cao hoặc thấp hơn trung bình):

\begin{equation}
    \text{sim}(u, v) = \frac{\displaystyle\sum_{i}
        (r_{ui} - \bar{r}_u)(r_{vi} - \bar{r}_v)}
        {\|r_u - \bar{r}_u\| \cdot \|r_v - \bar{r}_v\|}
    \label{eq:user_sim}
\end{equation}

Công thức dự đoán \textit{rating} có hiệu chỉnh \textit{bias}:

\begin{equation}
    \hat{r}_{ui} = \bar{r}_u +
        \frac{\displaystyle\sum_{v \in N(u)} \text{sim}(u,v)
              \cdot (r_{vi} - \bar{r}_v)}
             {\displaystyle\sum_{v \in N(u)} |\text{sim}(u,v)|}
    \label{eq:user_pred}
\end{equation}

Trong đó $\bar{r}_u$ là \textit{rating} trung bình của user $u$, $N(u)$ là tập $K$
láng giềng gần nhất của $u$ đã đánh giá sản phẩm $i$. Kết quả được giới hạn trong
khoảng $[1, 5]$ bằng \texttt{np.clip}. Toàn bộ ma trận \textit{similarity}
($943 \times 943$) được tính trước trong bước \texttt{fit()} bằng phép nhân ma trận
NumPy ($R_n \cdot R_n^T$), đảm bảo hiệu năng khi phục vụ nhiều yêu cầu gợi ý.

\subsection{Item-based Collaborative Filtering}
\label{section:2.2.3}

\textit{Item-based CF} tính độ tương đồng giữa các sản phẩm thay vì giữa người
dùng~\cite{sarwar2001item}. Phương pháp sử dụng \textbf{adjusted cosine similarity}
— trừ \textit{rating} trung bình của người dùng (không phải của sản phẩm) để loại
bỏ \textit{bias} đánh giá cá nhân khi so sánh hai sản phẩm:

\begin{equation}
    \text{sim}(i, j) = \frac{\displaystyle\sum_{u}
        (r_{ui} - \bar{r}_u)(r_{uj} - \bar{r}_u)}
        {\|r_i - \bar{r}\| \cdot \|r_j - \bar{r}\|}
    \label{eq:item_sim}
\end{equation}

Công thức dự đoán \textit{rating} cho user $u$ và item $i$, dựa trên $K = 30$ sản
phẩm tương đồng nhất mà user $u$ đã đánh giá:

\begin{equation}
    \hat{r}_{ui} =
        \frac{\displaystyle\sum_{j \in N(i,u)} \text{sim}(i,j) \cdot r_{uj}}
             {\displaystyle\sum_{j \in N(i,u)} |\text{sim}(i,j)|}
    \label{eq:item_pred}
\end{equation}

Công thức này sử dụng trực tiếp \textit{rating} thô $r_{uj}$ vì
\textit{adjusted cosine similarity} đã hấp thụ \textit{bias} người dùng vào bước
tính \textit{similarity}. Ma trận \textit{similarity} item ($1.682 \times 1.682$)
cũng được tính trước trong \texttt{fit()} để phục vụ dự đoán \textit{online} nhanh.

\textbf{Ưu điểm so với User-based CF:} Ma trận tương đồng item ổn định hơn theo thời
gian (danh sách phim ít thay đổi hơn người dùng) và có thể tính toán
\textit{offline}, phục vụ dự đoán \textit{online} nhanh hơn.

% ─────────────────────────────────────────────────────────
\section{Matrix Factorization}
\label{section:2.3}

\subsection{Ý tưởng cơ bản}
\label{section:2.3.1}

\textit{Matrix Factorization} (MF) là phương pháp thuộc nhóm \textit{Model-based CF},
nổi bật qua cuộc thi Netflix Prize (2009)~\cite{koren2009matrix}. Ý tưởng cốt lõi là
phân rã ma trận \textit{rating} $R$ ($n \times m$) thành tích của hai ma trận
\textit{latent factor} có chiều thấp hơn:

\begin{equation}
    R \approx P \times Q^T
    \label{eq:mf}
\end{equation}

Trong đó: $P$ ($n \times K$) là ma trận \textit{latent factor} của người dùng;
$Q$ ($m \times K$) là ma trận \textit{latent factor} của sản phẩm; $K$ là số
\textit{latent factors} (trong đồ án $K = 20$). \textit{Rating} được dự đoán bởi
tích vô hướng $\hat{r}_{ui} = p_u \cdot q_i$.

Các \textit{latent factor} không có nhãn tường minh, nhưng ngầm học các đặc trưng
ẩn của dữ liệu — ví dụ: thể loại phim, phong cách đạo diễn, mức độ kịch tính.
Người dùng và phim được ánh xạ vào cùng một không gian $K$ chiều, và
\textit{rating} được ước lượng qua mức độ ``khớp'' giữa hai vector.

Mô hình \textbf{Funk SVD}~\cite{funk2006netflix} mở rộng thêm \textit{bias} để tăng
độ chính xác:

\begin{equation}
    \hat{r}_{ui} = \mu + b_u + b_i + p_u \cdot q_i
    \label{eq:funk_svd}
\end{equation}

Trong đó $\mu$ là \textit{rating} trung bình toàn cục, $b_u$ là \textit{bias} của
user $u$ (xu hướng cho điểm cao/thấp hơn trung bình), $b_i$ là \textit{bias} của
item $i$ (phim được yêu thích hay không nói chung).

\subsection{Hàm mất mát và Regularization}
\label{section:2.3.2}

Mô hình tối thiểu hóa hàm mất mát bình phương có thêm \textit{regularization} L2
để tránh \textit{overfitting}:

\begin{equation}
    L = \sum_{(u,i) \in \text{rated}} (r_{ui} - \hat{r}_{ui})^2
      + \lambda \cdot \bigl(\|p_u\|^2 + \|q_i\|^2 + b_u^2 + b_i^2\bigr)
    \label{eq:loss}
\end{equation}

\textit{Regularization} với hệ số $\lambda = 0{,}02$ phạt các giá trị tham số quá
lớn, giúp mô hình tổng quát hóa tốt hơn trên dữ liệu chưa thấy. Không có
\textit{regularization}, mô hình dễ ghi nhớ tập \textit{train} mà không dự đoán
được tập \textit{test}.

\subsection{Stochastic Gradient Descent và theo dõi hội tụ}
\label{section:2.3.3}

Thuật toán \textit{Stochastic Gradient Descent} (SGD) cập nhật tham số sau mỗi
\textit{rating} đơn lẻ (thay vì toàn bộ dataset như Batch GD), giúp hội tụ nhanh
hơn~\cite{koren2009matrix}. Với \textit{residual}
$e_{ui} = r_{ui} - \hat{r}_{ui}$ và \textit{learning rate} $\alpha = 0{,}005$,
quy tắc cập nhật:

\begin{align}
    p_u &\leftarrow p_u + \alpha \cdot (e_{ui} \cdot q_i - \lambda \cdot p_u)
        \label{eq:sgd_p} \\
    q_i &\leftarrow q_i + \alpha \cdot (e_{ui} \cdot p_u - \lambda \cdot q_i)
        \label{eq:sgd_q} \\
    b_u &\leftarrow b_u + \alpha \cdot (e_{ui} - \lambda \cdot b_u)
        \label{eq:sgd_bu} \\
    b_i &\leftarrow b_i + \alpha \cdot (e_{ui} - \lambda \cdot b_i)
        \label{eq:sgd_bi}
\end{align}

Dữ liệu được xáo trộn ngẫu nhiên (\texttt{np.random.shuffle}) đầu mỗi \textit{epoch}
để tránh \textit{bias} thứ tự. Ma trận $P$ và $Q$ được khởi tạo từ phân phối chuẩn
$\mathcal{N}(0,\, 0{,}1)$; \textit{bias} $b_u$, $b_i$ khởi tạo bằng $0$. Mô hình
chạy $30$ \textit{epochs}. Sau mỗi \textit{epoch}, RMSE trên cả tập \textit{train}
và tập \textit{validation} được ghi lại để theo dõi hội tụ và phát hiện
\textit{overfitting}; hai đường \textit{loss curve} này được hiển thị trực quan
trên giao diện web \textbf{CineRec}.

% ─────────────────────────────────────────────────────────
\section{Các chỉ số đánh giá}
\label{section:2.4}

\subsection{Chiến lược chia dữ liệu Train/Test}
\label{section:2.4.1}

Dữ liệu được chia theo chiến lược \textbf{temporal split}: với mỗi người dùng,
$20\%$ \textit{rating} có \textit{timestamp} muộn nhất được giữ làm tập
\textit{test}, $80\%$ còn lại làm tập \textit{train}. Chiến lược này phản ánh đúng
kịch bản thực tế — mô hình chỉ học từ lịch sử quá khứ và dự đoán hành vi tương
lai — tránh hiện tượng \textit{data leakage} mà \textit{random split} có thể gây ra
(mô hình vô tình ``nhìn thấy tương lai'' khi huấn luyện).

\subsection{MAE và RMSE}
\label{section:2.4.2}

Hai chỉ số đánh giá chất lượng dự đoán \textit{rating} trên tập
\textit{test}~\cite{herlocker2004evaluating}:

\begin{equation}
    \text{MAE} = \frac{1}{n} \sum_{(u,i)} |r_{ui} - \hat{r}_{ui}|
    \label{eq:mae}
\end{equation}

\begin{equation}
    \text{RMSE} = \sqrt{\frac{1}{n} \sum_{(u,i)} (r_{ui} - \hat{r}_{ui})^2}
    \label{eq:rmse}
\end{equation}

MAE đo sai số tuyệt đối trung bình, dễ hiểu và ít nhạy với \textit{outlier}.
RMSE phạt nặng hơn các sai số lớn do bình phương hóa. Cả hai giá trị càng nhỏ,
mô hình dự đoán càng chính xác.

\subsection{Precision@K và Recall@K}
\label{section:2.4.3}

Hai chỉ số đánh giá chất lượng danh sách gợi ý Top-$K$, sử dụng ngưỡng
\textit{rating} $\geq 4{,}0$ để xác định sản phẩm ``relevant''
(khớp với tham số \texttt{threshold=4.0} trong hàm
\texttt{precision\_recall\_at\_k()}):

\begin{equation}
    \text{Precision@}K =
        \frac{|\{\text{relevant items}\} \cap \{\text{top-}K\text{ recs}\}|}{K}
    \label{eq:precision}
\end{equation}

\begin{equation}
    \text{Recall@}K =
        \frac{|\{\text{relevant items}\} \cap \{\text{top-}K\text{ recs}\}|}
             {|\{\text{relevant items}\}|}
    \label{eq:recall}
\end{equation}

$\text{Precision@}K$ đo tỷ lệ gợi ý chính xác trong Top-$K$;
$\text{Recall@}K$ đo tỷ lệ sản phẩm \textit{relevant} được tìm thấy trong Top-$K$.
Hai chỉ số này bổ sung cho MAE/RMSE vì đánh giá trực tiếp chất lượng danh sách gợi
ý thay vì độ chính xác từng \textit{rating} đơn lẻ. Do chi phí tính toán,
$\text{Precision@10}$ và $\text{Recall@10}$ được tính trên \textit{subset} $200$
người dùng đầu tiên của tập \textit{test}.

% ─────────────────────────────────────────────────────────
\section{Bộ dữ liệu MovieLens ml-100k}
\label{section:2.5}

\texttt{MovieLens ml-100k} là bộ dữ liệu chuẩn (\textit{benchmark}) được phát hành
bởi nhóm GroupLens Research, Đại học Minnesota~\cite{harper2015movielens}. Đây là
một trong những bộ dữ liệu được sử dụng rộng rãi nhất trong nghiên cứu hệ thống
gợi ý nhờ kích thước vừa phải, chất lượng tốt và tính khả dụng công khai.

\begin{table}[H]
    \centering
    \caption{Thống kê tổng quan dataset MovieLens ml-100k}
    \label{tab:ml100k}
    \begin{tabular}{lll}
        \toprule
        \textbf{Thuộc tính} & \textbf{Giá trị} & \textbf{Ghi chú} \\
        \midrule
        Số người dùng     & 943              & Đã ẩn danh \\
        Số phim           & 1.682            & Nhiều thể loại \\
        Tổng số rating    & 100.000          & Thang điểm 1--5 (nguyên) \\
        Rating trung bình & 3,53             & Lệch về phía cao \\
        Sparsity          & $\approx 93{,}70\%$
                          & $1 - \dfrac{100.000}{943 \times 1.682}$,
                            tính trên toàn bộ ma trận trước khi chia train/test \\
        Thời gian         & 1997--1998       & Unix timestamp \\
        \bottomrule
    \end{tabular}
\end{table}

Ma trận \textit{user-item} có độ thưa thớt cao ($\approx 93{,}70\%$), nghĩa là
trung bình mỗi người dùng chỉ đánh giá khoảng $6{,}3\%$ số phim có trong hệ thống.
Đây là thách thức điển hình của \textit{Collaborative Filtering}, đòi hỏi thuật
toán phải xử lý tốt với dữ liệu \textit{sparse}.

% ─────────────────────────────────────────────────────────
\section{Công nghệ và phương pháp triển khai}
\label{section:2.6}

\subsection{Triển khai thuật toán}
\label{section:2.6.1}

Toàn bộ ba mô hình \textit{Collaborative Filtering} (User-based CF, Item-based CF,
Matrix Factorization) được \textbf{triển khai hoàn toàn từ đầu} (\textit{from
scratch}) bằng NumPy, không sử dụng thư viện học máy chuyên dụng như
\texttt{Surprise}, \texttt{LightFM} hay các mô hình có sẵn của \texttt{scikit-learn}.
Điều này đảm bảo:

\begin{itemize}[itemsep=0.35em]
    \item Hiểu sâu cơ chế hoạt động của từng thuật toán (công thức tính
          \textit{similarity}, quy tắc dự đoán, cập nhật tham số SGD).
    \item Kiểm soát hoàn toàn luồng xử lý dữ liệu, \textit{hyperparameter tuning}
          và tối ưu hóa hiệu năng.
    \item Tính toán ma trận tương đồng bằng phép nhân ma trận vectorized
          (NumPy matrix multiplication), thay thế vòng lặp Python thuần, giảm thời
          gian tính từ phút xuống giây.
\end{itemize}

\subsection{Stack công nghệ}
\label{section:2.6.2}

\begin{table}[H]
    \centering
    \caption{Các công nghệ sử dụng trong hệ thống CineRec}
    \label{tab:tech}
    \begin{tabular}{l l p{6.8cm}}
        \toprule
        \textbf{Công nghệ} & \textbf{Phiên bản} & \textbf{Mục đích sử dụng} \\
        \midrule
        \texttt{Python}       & 3.10+  & Ngôn ngữ lập trình chính \\
        \texttt{NumPy}        & 1.24+  & \textbf{Lõi thuật toán}: Tính toán ma trận
                                \textit{similarity}, phép nhân ma trận
                                ($R_n \cdot R_n^T$), SGD cập nhật tham số
                                $P$, $Q$, $b_u$, $b_i$ \\
        \texttt{Pandas}       & 2.0+   & Đọc dữ liệu CSV, \textit{pivot table}
                                (xây dựng ma trận $R$), \textit{temporal split},
                                phân tích EDA \\
        \texttt{Scikit-learn} & 1.3+   & \textbf{Chỉ dùng cho đánh giá}:
                                \texttt{mean\_absolute\_error()},
                                \texttt{mean\_squared\_error()} để tính
                                MAE/RMSE. \textbf{Không} sử dụng cho xây
                                dựng mô hình CF hay MF \\
        \texttt{FastAPI}      & 0.110+ & Web framework backend RESTful API (endpoint
                                \texttt{/api/recommend}, \texttt{/api/stats}) \\
        \texttt{Uvicorn}      & 0.27+  & ASGI server chạy FastAPI (port 8000) \\
        \texttt{HTML5/CSS3/JS} & ---   & Giao diện web frontend SPA
                                (không dùng React/Vue, thuần Vanilla JS) \\
        \bottomrule
    \end{tabular}
\end{table}

\textbf{Tóm tắt kiến trúc:}
\begin{itemize}[itemsep=0.3em]
    \item \textbf{Data Layer:} Pandas đọc \texttt{ratings.csv}, \texttt{movies.csv};
          \texttt{temporal\_split()} chia train/test.
    \item \textbf{Model Layer:} Ba lớp \texttt{UserBasedCF}, \texttt{ItemBasedCF},
          \texttt{MatrixFactorization} cài đặt từ đầu bằng NumPy.
    \item \textbf{API Layer:} FastAPI expose endpoint RESTful, trả JSON; Uvicorn
          phục vụ HTTP.
    \item \textbf{Presentation Layer:} HTML/CSS/JS render giao diện, gọi API qua
          \texttt{fetch()}, vẽ biểu đồ \textit{loss curve} bằng Canvas API.
\end{itemize}

\end{document}
